\section{Conclusion}
\label{sec:conclusion}

We studied how far agreement among independent LLMs predicts the correctness of their answers.
We used 8 open-weight models from 7 families on factual, math and ethics tasks, plus a format-matched binary factual control.
Agreement is a strong label-free correctness signal for factual (AUROC 0.78--0.85) and math answers (0.85--0.91), and a much weaker one for ethical judgements (0.63--0.71).
Only part of the ethics gap reflects shared model bias.
The rest comes from the binary answer format and from noise in the crowd labels.
Plurality voting beats the best single model only on math.
It is neutral to slightly harmful on factual QA, and significantly harmful on ethics, where the deficit grows with every model added, as correlated errors would predict.
The key takeaway is that model consensus is a useful filter for facts and math, but not a substitute for human judgement on moral questions.

Several directions follow from this work.
First, repeating the study with frontier API models would test whether stronger models are also more correlated.
Second, chance-corrected agreement and binary versions of open-ended tasks would further separate format from domain.
Third, moral-judgement datasets with several annotations per item, such as Scruples~\citep{lourie2021scruples} or MoralChoice~\citep{scherrer2023moralchoice}, would let us test whether model agreement tracks {\em human} agreement rather than majority labels.
Finally, accuracy-weighted voting, agreement combined with self-probability, and a comparison with within-model self-consistency at matched compute~\citep{wang2023selfconsistency} could make consensus more reliable.
